\begin{table}[h]
  \centering
  \caption{Список значений и адресов операций с памятью для
  программы~\ref{fig:example-prog}}
  \small
  \begin{tabularx}{\textwidth}{@{}l X l@{}}
    \toprule
    Операций с памятью & \makecell[l]{Значение, записанное в
    память\\(только для W операций)} & Адрес ячейки памяти \\
    \midrule
    $W03$ & 0xdee7 & \([y, y+2)\) \\
    $W00$ & 0x18ae44139a7f23f9 & \([x, x+8)\) \\
    $R04$ &  & \([y, y+4)\) \\
    $W05$ & 0x6ce6ace0ddf15b17 & \([y, y+8)\) \\
    $W01$ & 0xcb5f722e & \([x, x+4)\) \\
    $W06$ & 0xff07b9b7354a2309 & \([y, y+8)\) \\
    $R02$ &  & \([x, x+8)\) \\
    $R07$ &  & \([y, y+1)\) \\
    $R13$ &  & \([z, z+2)\) \\
    $R10$ &  & \([x, x+8]\) \\
    $R14$ &  & \([z, z+2)\) \\
    $W15$ & 0x8a89de30 & \([z, z+4)\) \\
    $W12$ & 0x95 & \([x+6, x+7)\) \\
    $W16$ & 0xf7 & \([z+1, z+2)\) \\
    $W17$ & 0x1983 & \([z, z+2)\) \\
    $R23$ &  & \([w, w+1)\) \\
    $W20$ & 0x20E648DC & \([x+4, x+8)\) \\
    $R24$ &  & \([w, w+1)\) \\
    $R22$ &  & \([x+2, x+4)\) \\
    $W25$ & 0x523fbdbb & \([w, w+4)\) \\
    $R26$ &  & \([w+2, w+4)\) \\
    $R27$ &  & \([w+2, w+3)\) \\
    \bottomrule
  \end{tabularx}
  \label{tab:memory-access}
\end{table}
